% v23 research note: retain reusable equations, scope, and next frontier.
\subsection{通常型の全排除：外側 pair を早く使う}
\label{notes:g51523-ordinary}

\paragraph{確定した新しい結論．}
公開補節 \nolinkurl{g5_15_ordinary_outer_pair_obstruction_v23.tex} は，
標準五次整数接着の $G_{5,15}=K_5-\{23,24\}$ 通常型を全十四配置・全分母で排除する．
通常型の共有分母 $Q$ を増やす探索，非共鳴 $T\ge2$ の通常型 lifting，
通常型の外側四高 core 完成は今後の課題から外す．
第22版の共鳴・$T=1$ 排除は正しい局所補題として保持できるが，通常型の新しい証明には不要である．
第21版の正例外排除と，次節の負例外 $Q=3,5$ の全尺度排除を合わせ，
$G_{5,15}$ 全体が解決した．通常型の専用探索を再開しない．

\paragraph{再利用する三角形の一次式．}
中心と隣接三高 core の間の低辺を $2p$ とする．外側 core $i$ について
\[
 a=q_{2i},\quad c=p^2_{ip},\quad
 p^p_{2i}=-j,\quad p^p_{i2}=-\sigma,\quad q_{ip}=1+j\sigma
\]
と置く．node 変数 $u=1-n_{2i;p}$，$v=1-n_{ip;2}$ はともに $-4,\ldots,1$ であり，
$E=au-c$ とすると
\[
 \{E(jv+1)-aj\}\sigma=a-1-Ev,
 \qquad
 (1+j\sigma)\{jEv+E-aj\}=E-j.
\]
これは未完成の外側 core 内の一つの pair equation そのものである．
係数と定数がともに零となる共鳴は $E=j$，$v=(a-1)/j$ に限られる．
従って $a-1>j$ なら共鳴はなく，$q_{ip}\mid E-j\ne0$ から一様有限化できる．
中心の全体を完成してから外側 core を調べる順序は必須ではない．
他の支持でも，既知の低辺の両端から外側 core の必要 pair を先に課す方法を試す価値がある．
ただし $a-1\le j$ の共鳴は一般に除いてはならない．

\paragraph{通常型での有限表と最後の矛盾．}
外側二 core の尺度集合は，正側の順序を保って次の通りである：
\[
\begin{array}{c|c|c}
 (\text{中心},R^p,d)&S&T\\\hline
 (M_1,(-1,-4,4),5)&2&5\\
 (M_1,(-1,-1,1),-9)&2&\varnothing\\
 (M_2,(-1,-2,2),2)&2&1\\
 (M_2,(-2,-1,2),2)&\varnothing&4\\
 (M_3,(-1,-2,2),-7)&\varnothing&1\\
 (M_4,(-1,-2,2),9)&\varnothing&13
\end{array}
\]
第三行は $U=S+T-d-1=0$ で不可能．第一行だけが $U=1$ を残すが，
高分母 $(3,21,5)$ では元の pair・行和から
$HK=62$，$(H+1)M=14$，$M+N=6$ を得る．
最初の二つの整除性は $H=1,M=7$ を強制し，$N>0$ と矛盾する．
この表の全探索範囲は各外側入力の三十六 node 組から導かれ，任意の分母 cutoff ではない．
二つの負成分を交換した $M_2$ の三配置も独立に含めている．

\paragraph{局所分類との区別．}
通常型が消えたことから，五つの固定 $(R,D_0)$ の孤立局所問題がすべて不存在になったとはいえない．
例えば $R=(-1,-2,2)$，$D_0=2$，$(S,T,U)=(6,3,6)$ の正整数局所模型は実在し，
中心と外側 pair の組合せに接着できない．
独立した局所分類は他の支持に役立つ場合に限り続ける．
Schur の $k\ge0$ と $R=(-1,-m,m)$ の順を同時に固定しないという第22版の注意も引き継ぐ．
今回の証明は Schur の向きを使用していない．

\paragraph{検算の入口．}
\nolinkurl{verify_g5_15_ordinary_outer_pair_v23.py} は旧版 import なしで
四中心，十四配置，六 profile，十一個の相異なる外側入力を復元する．
node 直接列挙と約数列挙を独立に照合し，元の外側 pair equation を各残候補で確認する．
最後の一点は正約数の全検査と，二次残余 $2H^2-35H+124$ の判別式 $233$ の非平方性の
両方で確認する．後者は独立な補助確認であり，主証明は短い整除矛盾で完結する．
証明・ノート・検証器の hash と検証範囲を結果 JSON に記録する．
